% =====================================================================
%  BUỔI 7 — ĐÁP ÁN ĐỀ CHỮA (từ hoc-sinh/02-dap-an.md)
%  Sinh tự động bằng scripts/md2tex-chuade.py — không sửa tay.
% =====================================================================
\documentclass[11pt, a4paper]{article}
% =====================================================================
%  CAE SHANGHAI — GÓI LỆNH DÙNG CHUNG CHO TÀI LIỆU LUYỆN THI CSCA TOÁN
%  Biên dịch bằng XeLaTeX. Mọi file .tex trong buổi học \input file này.
% =====================================================================

% ---------- NGÔN NGỮ & PHÔNG CHỮ ----------
\usepackage{fontspec}
\usepackage{polyglossia}
\setmainlanguage{vietnamese}
\setotherlanguage{english}

% Chữ Latin (Anh + Việt): TeX Gyre Termes — bản clone Times New Roman,
% đi kèm TeX Live, hỗ trợ đầy đủ dấu tiếng Việt.
\setmainfont{TeX Gyre Termes}
\setsansfont{TeX Gyre Heros}

% Chữ Hán: ưu tiên Source Han Serif SC (思源宋体) — chuẩn học thuật Trung Quốc.
% Nếu máy chưa cài, tự động dùng Songti SC (宋体, có sẵn trên macOS).
% Cả hai đều thuộc họ 宋体 nên diện mạo tài liệu không đổi.
\IfFontExistsTF{Source Han Serif SC}{
  \newcommand{\cjkfamilyname}{Source Han Serif SC}
}{
  \IfFontExistsTF{Songti SC}{
    \newcommand{\cjkfamilyname}{Songti SC}
  }{
    \newcommand{\cjkfamilyname}{Noto Serif CJK SC}
  }
}

% xeCJK tự chuyển font khi gặp ký tự Hán, không cần bọc \cjk{} thủ công.
% Đây là điều kiện bắt buộc: font Latin (TeX Gyre Termes) không có glyph chữ Hán,
% nếu không khai báo xeCJK thì mọi chữ Hán viết trần sẽ bị nuốt khỏi bản in.
\usepackage{xeCJK}
\setCJKfamilyfont{caehan}{\cjkfamilyname}[Scale=0.95]
\setCJKmainfont{\cjkfamilyname}[Scale=0.95]
\setCJKsansfont{\cjkfamilyname}[Scale=0.95]
\setCJKmonofont{\cjkfamilyname}[Scale=0.95]

% \cjk{} ép dùng font Hán trong ngữ cảnh cần thiết (ví dụ trong tiêu đề in đậm
% hoặc trong bảng có định dạng riêng). Dùng \CJKfamily chứ không phải tên font,
% vì tên font chỉ là văn bản — sẽ in ra chữ thay vì đổi font.
\newcommand{\cjk}[1]{{\CJKfamily{caehan}#1}}

% ---------- TOÁN HỌC ----------
\usepackage{amsmath, amssymb, amsthm, mathtools}
\usepackage{unicode-math}
\setmathfont{TeX Gyre Termes Math}

% ---------- BỐ CỤC TRANG ----------
% Lề phải chừa vùng an toàn của letterhead:
%   banner trên chiếm 0–19mm  → top tối thiểu 2,2cm
%   banner dưới chiếm 281–297mm → bottom tối thiểu 2,2cm
% Lề trên 2,6cm (26mm) để tiêu đề đầu trang không đè banner (banner hết ở 19mm).
\usepackage[a4paper, top=2.6cm, bottom=2.4cm, left=2.5cm, right=2.0cm]{geometry}
\usepackage{fancyhdr}
\usepackage{titlesec}
\usepackage{enumitem}
\usepackage{graphicx}
\usepackage{booktabs, array, longtable, tabularx}
\usepackage{xcolor}
\usepackage[most]{tcolorbox}
\usepackage[colorlinks=true, linkcolor=cae_blue, citecolor=cae_blue, urlcolor=cae_blue]{hyperref}

% ---------- BẢNG MÀU THƯƠNG HIỆU ----------
\definecolor{cae_blue}{RGB}{0, 51, 102}
\definecolor{cae_gold}{RGB}{196, 149, 60}
\definecolor{box_bg}{RGB}{245, 247, 250}
\definecolor{box_line}{RGB}{0, 51, 102}

% ---------- TIÊU ĐỀ MỤC ----------
% Đệm đầu trang: banner letterhead chiếm 0–19mm, mà lề trên là 22mm.
% \section* đầu trang (sau \clearpage) sát lề nên phải chừa thêm, nếu không
% sẽ đè lên banner.
\titleformat{\section}
  {\normalfont\large\bfseries\sffamily\color{cae_blue}}{\thesection.}{0.6em}{}
\titleformat{\subsection}
  {\normalfont\normalsize\bfseries\sffamily\color{cae_blue!85!black}}{\thesubsection}{0.6em}{}
\titlespacing*{\section}{0pt}{14pt}{6pt}
\titlespacing*{\subsection}{0pt}{10pt}{4pt}

% Tiêu đề không đánh số ở đầu trang: \section* sau \clearpage nằm sát lề trên
% nên đè lên banner letterhead. Dùng \vspace* ngay TRƯỚC \section* không đủ,
% vì titlesec chèn khoảng cách riêng. Cách chắc chắn là đệm bằng \null.
\newcommand{\sectiondautrang}[1]{\section*{#1}}

% Dạng bài dùng \subsection*{} nên cần định dạng riêng, đậm và có gạch chân màu
\titleformat{\subsubsection}
  {\normalfont\normalsize\bfseries\sffamily\color{cae_blue!85!black}}{\thesubsubsection}{0.6em}{}

% ---------- LETTERHEAD (ảnh nền toàn trang) ----------
% Ảnh letterhead của CAE Shanghai: banner xanh trên/dưới kèm logo chìm giữa.
% Tỉ lệ ảnh 1414x2000 = 0.7070, khớp A4 (0.7071) nên chèn thẳng không méo.
% Đặt ở lớp nền dưới cùng để chữ không bị ảnh che.
% Dùng \AddToShipoutPictureBG (không có dấu *) để áp dụng cho MỌI trang.
\usepackage{eso-pic}
\AddToShipoutPictureBG{%
  \AtPageLowerLeft{%
    \includegraphics[width=\paperwidth,height=\paperheight]{images/letterhead.png}%
  }%
}

% ---------- ĐẦU / CHÂN TRANG ----------
% Letterhead đã có banner trên và dưới, nên bỏ dòng chạy để tránh chồng chữ.
% Số trang đặt trong vùng nội dung, sát trên banner dưới, lề phải.
\pagestyle{fancy}
\setlength{\headheight}{0pt}
\fancyhf{}
\renewcommand{\headrulewidth}{0pt}
\renewcommand{\footrulewidth}{0pt}
\fancyfoot[R]{\raisebox{0.9cm}{\small\color{cae_blue}\thepage}}

% ---------- MÔI TRƯỜNG NỘI DUNG ----------
% Trong tcolorbox, LaTeX đặt parindent = parskip = 0, nên các đoạn văn trong
% khung dính liền nhau. Đặt lại khoảng cách đoạn cho mọi khung.
\newcommand{\khoangcachdoan}{\setlength{\parskip}{0.5em}\setlength{\parindent}{0pt}}
\tcbset{before upper={\khoangcachdoan}}

\newtcolorbox{dinhnghia}[1][]{
  colback=box_bg, colframe=box_line, boxrule=0.6pt, arc=1pt,
  fonttitle=\bfseries, title={Định nghĩa #1}, breakable}
\newtcolorbox{dinhly}[1][]{
  colback=box_bg, colframe=cae_gold, boxrule=0.6pt, arc=1pt,
  fonttitle=\bfseries, title={Tính chất #1}, breakable}
\newtcolorbox{vidu}[1][]{
  colback=white, colframe=black!40, boxrule=0.4pt, arc=1pt,
  fonttitle=\bfseries, title={Ví dụ #1}, breakable}
\newtcolorbox{baitap}[1][]{
  colback=white, colframe=cae_blue!50, boxrule=0.4pt, arc=1pt,
  fonttitle=\bfseries, title={Câu #1}, breakable}

% Ô đáp án: viền xanh đậm hơn, dùng trong file đáp án
\newtcolorbox{dapan}[1][]{
  colback=box_bg, colframe=cae_blue, boxrule=0.6pt, arc=1pt,
  fonttitle=\bfseries, title={Câu #1}, breakable}

% ---------- TÊN CHỨNG MINH ----------
\renewcommand{\proofname}{\normalfont\bfseries Chứng minh}

% ---------- MÔN HỌC ----------
% Đổi được ở file .tex trước khi % =====================================================================
%  CAE SHANGHAI — GÓI LỆNH DÙNG CHUNG CHO TÀI LIỆU LUYỆN THI CSCA TOÁN
%  Biên dịch bằng XeLaTeX. Mọi file .tex trong buổi học \input file này.
% =====================================================================

% ---------- NGÔN NGỮ & PHÔNG CHỮ ----------
\usepackage{fontspec}
\usepackage{polyglossia}
\setmainlanguage{vietnamese}
\setotherlanguage{english}

% Chữ Latin (Anh + Việt): TeX Gyre Termes — bản clone Times New Roman,
% đi kèm TeX Live, hỗ trợ đầy đủ dấu tiếng Việt.
\setmainfont{TeX Gyre Termes}
\setsansfont{TeX Gyre Heros}

% Chữ Hán: ưu tiên Source Han Serif SC (思源宋体) — chuẩn học thuật Trung Quốc.
% Nếu máy chưa cài, tự động dùng Songti SC (宋体, có sẵn trên macOS).
% Cả hai đều thuộc họ 宋体 nên diện mạo tài liệu không đổi.
\IfFontExistsTF{Source Han Serif SC}{
  \newcommand{\cjkfamilyname}{Source Han Serif SC}
}{
  \IfFontExistsTF{Songti SC}{
    \newcommand{\cjkfamilyname}{Songti SC}
  }{
    \newcommand{\cjkfamilyname}{Noto Serif CJK SC}
  }
}

% xeCJK tự chuyển font khi gặp ký tự Hán, không cần bọc \cjk{} thủ công.
% Đây là điều kiện bắt buộc: font Latin (TeX Gyre Termes) không có glyph chữ Hán,
% nếu không khai báo xeCJK thì mọi chữ Hán viết trần sẽ bị nuốt khỏi bản in.
\usepackage{xeCJK}
\setCJKfamilyfont{caehan}{\cjkfamilyname}[Scale=0.95]
\setCJKmainfont{\cjkfamilyname}[Scale=0.95]
\setCJKsansfont{\cjkfamilyname}[Scale=0.95]
\setCJKmonofont{\cjkfamilyname}[Scale=0.95]

% \cjk{} ép dùng font Hán trong ngữ cảnh cần thiết (ví dụ trong tiêu đề in đậm
% hoặc trong bảng có định dạng riêng). Dùng \CJKfamily chứ không phải tên font,
% vì tên font chỉ là văn bản — sẽ in ra chữ thay vì đổi font.
\newcommand{\cjk}[1]{{\CJKfamily{caehan}#1}}

% ---------- TOÁN HỌC ----------
\usepackage{amsmath, amssymb, amsthm, mathtools}
\usepackage{unicode-math}
\setmathfont{TeX Gyre Termes Math}

% ---------- BỐ CỤC TRANG ----------
% Lề phải chừa vùng an toàn của letterhead:
%   banner trên chiếm 0–19mm  → top tối thiểu 2,2cm
%   banner dưới chiếm 281–297mm → bottom tối thiểu 2,2cm
% Lề trên 2,6cm (26mm) để tiêu đề đầu trang không đè banner (banner hết ở 19mm).
\usepackage[a4paper, top=2.6cm, bottom=2.4cm, left=2.5cm, right=2.0cm]{geometry}
\usepackage{fancyhdr}
\usepackage{titlesec}
\usepackage{enumitem}
\usepackage{graphicx}
\usepackage{booktabs, array, longtable, tabularx}
\usepackage{xcolor}
\usepackage[most]{tcolorbox}
\usepackage[colorlinks=true, linkcolor=cae_blue, citecolor=cae_blue, urlcolor=cae_blue]{hyperref}

% ---------- BẢNG MÀU THƯƠNG HIỆU ----------
\definecolor{cae_blue}{RGB}{0, 51, 102}
\definecolor{cae_gold}{RGB}{196, 149, 60}
\definecolor{box_bg}{RGB}{245, 247, 250}
\definecolor{box_line}{RGB}{0, 51, 102}

% ---------- TIÊU ĐỀ MỤC ----------
% Đệm đầu trang: banner letterhead chiếm 0–19mm, mà lề trên là 22mm.
% \section* đầu trang (sau \clearpage) sát lề nên phải chừa thêm, nếu không
% sẽ đè lên banner.
\titleformat{\section}
  {\normalfont\large\bfseries\sffamily\color{cae_blue}}{\thesection.}{0.6em}{}
\titleformat{\subsection}
  {\normalfont\normalsize\bfseries\sffamily\color{cae_blue!85!black}}{\thesubsection}{0.6em}{}
\titlespacing*{\section}{0pt}{14pt}{6pt}
\titlespacing*{\subsection}{0pt}{10pt}{4pt}

% Tiêu đề không đánh số ở đầu trang: \section* sau \clearpage nằm sát lề trên
% nên đè lên banner letterhead. Dùng \vspace* ngay TRƯỚC \section* không đủ,
% vì titlesec chèn khoảng cách riêng. Cách chắc chắn là đệm bằng \null.
\newcommand{\sectiondautrang}[1]{\section*{#1}}

% Dạng bài dùng \subsection*{} nên cần định dạng riêng, đậm và có gạch chân màu
\titleformat{\subsubsection}
  {\normalfont\normalsize\bfseries\sffamily\color{cae_blue!85!black}}{\thesubsubsection}{0.6em}{}

% ---------- LETTERHEAD (ảnh nền toàn trang) ----------
% Ảnh letterhead của CAE Shanghai: banner xanh trên/dưới kèm logo chìm giữa.
% Tỉ lệ ảnh 1414x2000 = 0.7070, khớp A4 (0.7071) nên chèn thẳng không méo.
% Đặt ở lớp nền dưới cùng để chữ không bị ảnh che.
% Dùng \AddToShipoutPictureBG (không có dấu *) để áp dụng cho MỌI trang.
\usepackage{eso-pic}
\AddToShipoutPictureBG{%
  \AtPageLowerLeft{%
    \includegraphics[width=\paperwidth,height=\paperheight]{images/letterhead.png}%
  }%
}

% ---------- ĐẦU / CHÂN TRANG ----------
% Letterhead đã có banner trên và dưới, nên bỏ dòng chạy để tránh chồng chữ.
% Số trang đặt trong vùng nội dung, sát trên banner dưới, lề phải.
\pagestyle{fancy}
\setlength{\headheight}{0pt}
\fancyhf{}
\renewcommand{\headrulewidth}{0pt}
\renewcommand{\footrulewidth}{0pt}
\fancyfoot[R]{\raisebox{0.9cm}{\small\color{cae_blue}\thepage}}

% ---------- MÔI TRƯỜNG NỘI DUNG ----------
% Trong tcolorbox, LaTeX đặt parindent = parskip = 0, nên các đoạn văn trong
% khung dính liền nhau. Đặt lại khoảng cách đoạn cho mọi khung.
\newcommand{\khoangcachdoan}{\setlength{\parskip}{0.5em}\setlength{\parindent}{0pt}}
\tcbset{before upper={\khoangcachdoan}}

\newtcolorbox{dinhnghia}[1][]{
  colback=box_bg, colframe=box_line, boxrule=0.6pt, arc=1pt,
  fonttitle=\bfseries, title={Định nghĩa #1}, breakable}
\newtcolorbox{dinhly}[1][]{
  colback=box_bg, colframe=cae_gold, boxrule=0.6pt, arc=1pt,
  fonttitle=\bfseries, title={Tính chất #1}, breakable}
\newtcolorbox{vidu}[1][]{
  colback=white, colframe=black!40, boxrule=0.4pt, arc=1pt,
  fonttitle=\bfseries, title={Ví dụ #1}, breakable}
\newtcolorbox{baitap}[1][]{
  colback=white, colframe=cae_blue!50, boxrule=0.4pt, arc=1pt,
  fonttitle=\bfseries, title={Câu #1}, breakable}

% Ô đáp án: viền xanh đậm hơn, dùng trong file đáp án
\newtcolorbox{dapan}[1][]{
  colback=box_bg, colframe=cae_blue, boxrule=0.6pt, arc=1pt,
  fonttitle=\bfseries, title={Câu #1}, breakable}

% ---------- TÊN CHỨNG MINH ----------
\renewcommand{\proofname}{\normalfont\bfseries Chứng minh}

% ---------- MÔN HỌC ----------
% Đổi được ở file .tex trước khi % =====================================================================
%  CAE SHANGHAI — GÓI LỆNH DÙNG CHUNG CHO TÀI LIỆU LUYỆN THI CSCA TOÁN
%  Biên dịch bằng XeLaTeX. Mọi file .tex trong buổi học \input file này.
% =====================================================================

% ---------- NGÔN NGỮ & PHÔNG CHỮ ----------
\usepackage{fontspec}
\usepackage{polyglossia}
\setmainlanguage{vietnamese}
\setotherlanguage{english}

% Chữ Latin (Anh + Việt): TeX Gyre Termes — bản clone Times New Roman,
% đi kèm TeX Live, hỗ trợ đầy đủ dấu tiếng Việt.
\setmainfont{TeX Gyre Termes}
\setsansfont{TeX Gyre Heros}

% Chữ Hán: ưu tiên Source Han Serif SC (思源宋体) — chuẩn học thuật Trung Quốc.
% Nếu máy chưa cài, tự động dùng Songti SC (宋体, có sẵn trên macOS).
% Cả hai đều thuộc họ 宋体 nên diện mạo tài liệu không đổi.
\IfFontExistsTF{Source Han Serif SC}{
  \newcommand{\cjkfamilyname}{Source Han Serif SC}
}{
  \IfFontExistsTF{Songti SC}{
    \newcommand{\cjkfamilyname}{Songti SC}
  }{
    \newcommand{\cjkfamilyname}{Noto Serif CJK SC}
  }
}

% xeCJK tự chuyển font khi gặp ký tự Hán, không cần bọc \cjk{} thủ công.
% Đây là điều kiện bắt buộc: font Latin (TeX Gyre Termes) không có glyph chữ Hán,
% nếu không khai báo xeCJK thì mọi chữ Hán viết trần sẽ bị nuốt khỏi bản in.
\usepackage{xeCJK}
\setCJKfamilyfont{caehan}{\cjkfamilyname}[Scale=0.95]
\setCJKmainfont{\cjkfamilyname}[Scale=0.95]
\setCJKsansfont{\cjkfamilyname}[Scale=0.95]
\setCJKmonofont{\cjkfamilyname}[Scale=0.95]

% \cjk{} ép dùng font Hán trong ngữ cảnh cần thiết (ví dụ trong tiêu đề in đậm
% hoặc trong bảng có định dạng riêng). Dùng \CJKfamily chứ không phải tên font,
% vì tên font chỉ là văn bản — sẽ in ra chữ thay vì đổi font.
\newcommand{\cjk}[1]{{\CJKfamily{caehan}#1}}

% ---------- TOÁN HỌC ----------
\usepackage{amsmath, amssymb, amsthm, mathtools}
\usepackage{unicode-math}
\setmathfont{TeX Gyre Termes Math}

% ---------- BỐ CỤC TRANG ----------
% Lề phải chừa vùng an toàn của letterhead:
%   banner trên chiếm 0–19mm  → top tối thiểu 2,2cm
%   banner dưới chiếm 281–297mm → bottom tối thiểu 2,2cm
% Lề trên 2,6cm (26mm) để tiêu đề đầu trang không đè banner (banner hết ở 19mm).
\usepackage[a4paper, top=2.6cm, bottom=2.4cm, left=2.5cm, right=2.0cm]{geometry}
\usepackage{fancyhdr}
\usepackage{titlesec}
\usepackage{enumitem}
\usepackage{graphicx}
\usepackage{booktabs, array, longtable, tabularx}
\usepackage{xcolor}
\usepackage[most]{tcolorbox}
\usepackage[colorlinks=true, linkcolor=cae_blue, citecolor=cae_blue, urlcolor=cae_blue]{hyperref}

% ---------- BẢNG MÀU THƯƠNG HIỆU ----------
\definecolor{cae_blue}{RGB}{0, 51, 102}
\definecolor{cae_gold}{RGB}{196, 149, 60}
\definecolor{box_bg}{RGB}{245, 247, 250}
\definecolor{box_line}{RGB}{0, 51, 102}

% ---------- TIÊU ĐỀ MỤC ----------
% Đệm đầu trang: banner letterhead chiếm 0–19mm, mà lề trên là 22mm.
% \section* đầu trang (sau \clearpage) sát lề nên phải chừa thêm, nếu không
% sẽ đè lên banner.
\titleformat{\section}
  {\normalfont\large\bfseries\sffamily\color{cae_blue}}{\thesection.}{0.6em}{}
\titleformat{\subsection}
  {\normalfont\normalsize\bfseries\sffamily\color{cae_blue!85!black}}{\thesubsection}{0.6em}{}
\titlespacing*{\section}{0pt}{14pt}{6pt}
\titlespacing*{\subsection}{0pt}{10pt}{4pt}

% Tiêu đề không đánh số ở đầu trang: \section* sau \clearpage nằm sát lề trên
% nên đè lên banner letterhead. Dùng \vspace* ngay TRƯỚC \section* không đủ,
% vì titlesec chèn khoảng cách riêng. Cách chắc chắn là đệm bằng \null.
\newcommand{\sectiondautrang}[1]{\section*{#1}}

% Dạng bài dùng \subsection*{} nên cần định dạng riêng, đậm và có gạch chân màu
\titleformat{\subsubsection}
  {\normalfont\normalsize\bfseries\sffamily\color{cae_blue!85!black}}{\thesubsubsection}{0.6em}{}

% ---------- LETTERHEAD (ảnh nền toàn trang) ----------
% Ảnh letterhead của CAE Shanghai: banner xanh trên/dưới kèm logo chìm giữa.
% Tỉ lệ ảnh 1414x2000 = 0.7070, khớp A4 (0.7071) nên chèn thẳng không méo.
% Đặt ở lớp nền dưới cùng để chữ không bị ảnh che.
% Dùng \AddToShipoutPictureBG (không có dấu *) để áp dụng cho MỌI trang.
\usepackage{eso-pic}
\AddToShipoutPictureBG{%
  \AtPageLowerLeft{%
    \includegraphics[width=\paperwidth,height=\paperheight]{images/letterhead.png}%
  }%
}

% ---------- ĐẦU / CHÂN TRANG ----------
% Letterhead đã có banner trên và dưới, nên bỏ dòng chạy để tránh chồng chữ.
% Số trang đặt trong vùng nội dung, sát trên banner dưới, lề phải.
\pagestyle{fancy}
\setlength{\headheight}{0pt}
\fancyhf{}
\renewcommand{\headrulewidth}{0pt}
\renewcommand{\footrulewidth}{0pt}
\fancyfoot[R]{\raisebox{0.9cm}{\small\color{cae_blue}\thepage}}

% ---------- MÔI TRƯỜNG NỘI DUNG ----------
% Trong tcolorbox, LaTeX đặt parindent = parskip = 0, nên các đoạn văn trong
% khung dính liền nhau. Đặt lại khoảng cách đoạn cho mọi khung.
\newcommand{\khoangcachdoan}{\setlength{\parskip}{0.5em}\setlength{\parindent}{0pt}}
\tcbset{before upper={\khoangcachdoan}}

\newtcolorbox{dinhnghia}[1][]{
  colback=box_bg, colframe=box_line, boxrule=0.6pt, arc=1pt,
  fonttitle=\bfseries, title={Định nghĩa #1}, breakable}
\newtcolorbox{dinhly}[1][]{
  colback=box_bg, colframe=cae_gold, boxrule=0.6pt, arc=1pt,
  fonttitle=\bfseries, title={Tính chất #1}, breakable}
\newtcolorbox{vidu}[1][]{
  colback=white, colframe=black!40, boxrule=0.4pt, arc=1pt,
  fonttitle=\bfseries, title={Ví dụ #1}, breakable}
\newtcolorbox{baitap}[1][]{
  colback=white, colframe=cae_blue!50, boxrule=0.4pt, arc=1pt,
  fonttitle=\bfseries, title={Câu #1}, breakable}

% Ô đáp án: viền xanh đậm hơn, dùng trong file đáp án
\newtcolorbox{dapan}[1][]{
  colback=box_bg, colframe=cae_blue, boxrule=0.6pt, arc=1pt,
  fonttitle=\bfseries, title={Câu #1}, breakable}

% ---------- TÊN CHỨNG MINH ----------
\renewcommand{\proofname}{\normalfont\bfseries Chứng minh}

% ---------- MÔN HỌC ----------
% Đổi được ở file .tex trước khi \input{cae-style} bằng \renewcommand{\monhoc}{...}
\providecommand{\monhoc}{Môn Toán}

% ---------- TRANG BÌA GỌN ----------
% #1 = loại tài liệu (TÀI LIỆU CHUẨN BỊ, TÀI LIỆU BUỔI HỌC, ...)
% #2 = số buổi
% #3 = tên buổi tiếng Việt
% #4 = tên buổi tiếng Hán
% #5 = module và tỉ trọng
\newcommand{\trangbia}[5]{%
  \begin{center}
    \vspace*{0.6cm}
    {\large\bfseries\color{cae_blue} CAE SHANGHAI}\\[0.15cm]
    {\footnotesize Tài liệu luyện thi CSCA · \monhoc}\\[0.5cm]
    {\color{cae_gold}\rule{0.55\linewidth}{0.8pt}}\\[0.35cm]
    {\small\bfseries\color{cae_blue!80!black} #1}\\[0.35cm]
    {\Large\bfseries\color{cae_blue} BUỔI #2 — #3}\\[0.2cm]
    {\large\color{cae_blue!85!black} \cjk{#4}}\\[0.25cm]
    \ifx\relax#5\relax\else{\small #5}\\[0.35cm]\fi
    {\color{cae_gold}\rule{0.55\linewidth}{0.8pt}}\\[0.3cm]
    {\footnotesize Biên soạn: Hoàng Dương, CAE Shanghai \quad·\quad \today}
  \end{center}
  \vspace{0.3cm}
}

% ---------- TRANG BÌA CHO ĐỀ KIỂM TRA ----------
% Khác \trangbia ở chỗ không có chữ "BUỔI N" — dùng cho đề thi, không phải bài giảng.
% #1 = loại tài liệu, #2 = tên kỳ thi (Việt), #3 = tên kỳ thi (Hán), #4 = thông số
\newcommand{\biade}[4]{%
  \begin{center}
    \vspace*{0.6cm}
    {\large\bfseries\color{cae_blue} CAE SHANGHAI}\\[0.15cm]
    {\footnotesize Tài liệu luyện thi CSCA · \monhoc}\\[0.5cm]
    {\color{cae_gold}\rule{0.55\linewidth}{0.8pt}}\\[0.35cm]
    {\small\bfseries\color{cae_blue!80!black} #1}\\[0.35cm]
    {\Large\bfseries\color{cae_blue} #2}\\[0.2cm]
    {\large\color{cae_blue!85!black} \cjk{#3}}\\[0.25cm]
    {\small #4}\\[0.35cm]
    {\color{cae_gold}\rule{0.55\linewidth}{0.8pt}}\\[0.3cm]
    {\footnotesize Biên soạn: Hoàng Dương, CAE Shanghai \quad·\quad \today}
  \end{center}
  \vspace{0.3cm}
}

% ---------- TIỆN ÍCH ----------
% Nhãn từ khóa in đậm ở đầu đoạn, dùng cho Từ vựng / Từ khóa nhận dạng / Chú ý
% Khoảng cách đoạn do khung tự đặt (xem \khoangcachdoan ở trên),
% ở đây chỉ cần tách đoạn và bỏ thụt đầu dòng.
\newcommand{\nhan}[1]{\par\noindent\textbf{#1.}\ \ignorespaces}

% Bảng gọn, không đánh số
\newcommand{\bang}[1]{%
  \begin{center}#1\end{center}}

% Đánh dấu ô trống cho học sinh điền
\newcommand{\chotrong}{\makebox[3.2cm]{\dotfill}}
 bằng \renewcommand{\monhoc}{...}
\providecommand{\monhoc}{Môn Toán}

% ---------- TRANG BÌA GỌN ----------
% #1 = loại tài liệu (TÀI LIỆU CHUẨN BỊ, TÀI LIỆU BUỔI HỌC, ...)
% #2 = số buổi
% #3 = tên buổi tiếng Việt
% #4 = tên buổi tiếng Hán
% #5 = module và tỉ trọng
\newcommand{\trangbia}[5]{%
  \begin{center}
    \vspace*{0.6cm}
    {\large\bfseries\color{cae_blue} CAE SHANGHAI}\\[0.15cm]
    {\footnotesize Tài liệu luyện thi CSCA · \monhoc}\\[0.5cm]
    {\color{cae_gold}\rule{0.55\linewidth}{0.8pt}}\\[0.35cm]
    {\small\bfseries\color{cae_blue!80!black} #1}\\[0.35cm]
    {\Large\bfseries\color{cae_blue} BUỔI #2 — #3}\\[0.2cm]
    {\large\color{cae_blue!85!black} \cjk{#4}}\\[0.25cm]
    \ifx\relax#5\relax\else{\small #5}\\[0.35cm]\fi
    {\color{cae_gold}\rule{0.55\linewidth}{0.8pt}}\\[0.3cm]
    {\footnotesize Biên soạn: Hoàng Dương, CAE Shanghai \quad·\quad \today}
  \end{center}
  \vspace{0.3cm}
}

% ---------- TRANG BÌA CHO ĐỀ KIỂM TRA ----------
% Khác \trangbia ở chỗ không có chữ "BUỔI N" — dùng cho đề thi, không phải bài giảng.
% #1 = loại tài liệu, #2 = tên kỳ thi (Việt), #3 = tên kỳ thi (Hán), #4 = thông số
\newcommand{\biade}[4]{%
  \begin{center}
    \vspace*{0.6cm}
    {\large\bfseries\color{cae_blue} CAE SHANGHAI}\\[0.15cm]
    {\footnotesize Tài liệu luyện thi CSCA · \monhoc}\\[0.5cm]
    {\color{cae_gold}\rule{0.55\linewidth}{0.8pt}}\\[0.35cm]
    {\small\bfseries\color{cae_blue!80!black} #1}\\[0.35cm]
    {\Large\bfseries\color{cae_blue} #2}\\[0.2cm]
    {\large\color{cae_blue!85!black} \cjk{#3}}\\[0.25cm]
    {\small #4}\\[0.35cm]
    {\color{cae_gold}\rule{0.55\linewidth}{0.8pt}}\\[0.3cm]
    {\footnotesize Biên soạn: Hoàng Dương, CAE Shanghai \quad·\quad \today}
  \end{center}
  \vspace{0.3cm}
}

% ---------- TIỆN ÍCH ----------
% Nhãn từ khóa in đậm ở đầu đoạn, dùng cho Từ vựng / Từ khóa nhận dạng / Chú ý
% Khoảng cách đoạn do khung tự đặt (xem \khoangcachdoan ở trên),
% ở đây chỉ cần tách đoạn và bỏ thụt đầu dòng.
\newcommand{\nhan}[1]{\par\noindent\textbf{#1.}\ \ignorespaces}

% Bảng gọn, không đánh số
\newcommand{\bang}[1]{%
  \begin{center}#1\end{center}}

% Đánh dấu ô trống cho học sinh điền
\newcommand{\chotrong}{\makebox[3.2cm]{\dotfill}}
 bằng \renewcommand{\monhoc}{...}
\providecommand{\monhoc}{Môn Toán}

% ---------- TRANG BÌA GỌN ----------
% #1 = loại tài liệu (TÀI LIỆU CHUẨN BỊ, TÀI LIỆU BUỔI HỌC, ...)
% #2 = số buổi
% #3 = tên buổi tiếng Việt
% #4 = tên buổi tiếng Hán
% #5 = module và tỉ trọng
\newcommand{\trangbia}[5]{%
  \begin{center}
    \vspace*{0.6cm}
    {\large\bfseries\color{cae_blue} CAE SHANGHAI}\\[0.15cm]
    {\footnotesize Tài liệu luyện thi CSCA · \monhoc}\\[0.5cm]
    {\color{cae_gold}\rule{0.55\linewidth}{0.8pt}}\\[0.35cm]
    {\small\bfseries\color{cae_blue!80!black} #1}\\[0.35cm]
    {\Large\bfseries\color{cae_blue} BUỔI #2 — #3}\\[0.2cm]
    {\large\color{cae_blue!85!black} \cjk{#4}}\\[0.25cm]
    \ifx\relax#5\relax\else{\small #5}\\[0.35cm]\fi
    {\color{cae_gold}\rule{0.55\linewidth}{0.8pt}}\\[0.3cm]
    {\footnotesize Biên soạn: Hoàng Dương, CAE Shanghai \quad·\quad \today}
  \end{center}
  \vspace{0.3cm}
}

% ---------- TRANG BÌA CHO ĐỀ KIỂM TRA ----------
% Khác \trangbia ở chỗ không có chữ "BUỔI N" — dùng cho đề thi, không phải bài giảng.
% #1 = loại tài liệu, #2 = tên kỳ thi (Việt), #3 = tên kỳ thi (Hán), #4 = thông số
\newcommand{\biade}[4]{%
  \begin{center}
    \vspace*{0.6cm}
    {\large\bfseries\color{cae_blue} CAE SHANGHAI}\\[0.15cm]
    {\footnotesize Tài liệu luyện thi CSCA · \monhoc}\\[0.5cm]
    {\color{cae_gold}\rule{0.55\linewidth}{0.8pt}}\\[0.35cm]
    {\small\bfseries\color{cae_blue!80!black} #1}\\[0.35cm]
    {\Large\bfseries\color{cae_blue} #2}\\[0.2cm]
    {\large\color{cae_blue!85!black} \cjk{#3}}\\[0.25cm]
    {\small #4}\\[0.35cm]
    {\color{cae_gold}\rule{0.55\linewidth}{0.8pt}}\\[0.3cm]
    {\footnotesize Biên soạn: Hoàng Dương, CAE Shanghai \quad·\quad \today}
  \end{center}
  \vspace{0.3cm}
}

% ---------- TIỆN ÍCH ----------
% Nhãn từ khóa in đậm ở đầu đoạn, dùng cho Từ vựng / Từ khóa nhận dạng / Chú ý
% Khoảng cách đoạn do khung tự đặt (xem \khoangcachdoan ở trên),
% ở đây chỉ cần tách đoạn và bỏ thụt đầu dòng.
\newcommand{\nhan}[1]{\par\noindent\textbf{#1.}\ \ignorespaces}

% Bảng gọn, không đánh số
\newcommand{\bang}[1]{%
  \begin{center}#1\end{center}}

% Đánh dấu ô trống cho học sinh điền
\newcommand{\chotrong}{\makebox[3.2cm]{\dotfill}}


\begin{document}

\trangbia{ĐÁP ÁN ĐỀ CHỮA}{7}{Tập hợp, hàm số, lượng giác, dãy số}{综合复习}{Đáp án và lời giải}

Mở file này \textbf{sau khi đã làm xong toàn bộ 48 câu trong 60 phút}.

\section{BẢNG ĐÁP ÁN}

\bang{
\begin{tabular}{@{}lcccccccccccc@{}}
\toprule
\textbf{Câu} & \textbf{1} & \textbf{2} & \textbf{3} & \textbf{4} & \textbf{5} & \textbf{6} & \textbf{7} & \textbf{8} & \textbf{9} & \textbf{10} & \textbf{11} & \textbf{12} \\
\midrule
Đáp án & C & B & C & A & A & B & C & A & C & C & D & D \\
\bottomrule
\end{tabular}}

\bang{
\begin{tabular}{@{}lcccccccccccc@{}}
\toprule
\textbf{Câu} & \textbf{13} & \textbf{14} & \textbf{15} & \textbf{16} & \textbf{17} & \textbf{18} & \textbf{19} & \textbf{20} & \textbf{21} & \textbf{22} & \textbf{23} & \textbf{24} \\
\midrule
Đáp án & B & B & C & D & B & B & B & A & D & C & D & C \\
\bottomrule
\end{tabular}}

\bang{
\begin{tabular}{@{}lcccccccccccc@{}}
\toprule
\textbf{Câu} & \textbf{25} & \textbf{26} & \textbf{27} & \textbf{28} & \textbf{29} & \textbf{30} & \textbf{31} & \textbf{32} & \textbf{33} & \textbf{34} & \textbf{35} & \textbf{36} \\
\midrule
Đáp án & A & D & B & D & A & A & A & A & C & B & D & C \\
\bottomrule
\end{tabular}}

\bang{
\begin{tabular}{@{}lcccccccccccc@{}}
\toprule
\textbf{Câu} & \textbf{37} & \textbf{38} & \textbf{39} & \textbf{40} & \textbf{41} & \textbf{42} & \textbf{43} & \textbf{44} & \textbf{45} & \textbf{46} & \textbf{47} & \textbf{48} \\
\midrule
Đáp án & A & C & B & D & C & D & D & A & B & B & D & A \\
\bottomrule
\end{tabular}}

\section{PHẦN I — TẬP HỢP VÀ BẤT ĐẲNG THỨC}

\begin{dapan}[1 — Đáp án C]

\textbf{Giải.} $A\subseteq B$ nên mọi phần tử của $A$ phải thuộc $B$. Phần tử $1$ đã có trong $B$. Còn $a$ phải là một phần tử của $B$ và khác $1$, vậy $a\in\{3,6,7,9\}$, có $4$ giá trị.


\end{dapan}
\vspace{0.3cm}
\begin{dapan}[2 — Đáp án B]

\textbf{Giải.} $x^2-4<0 \iff -2<x<2$, vậy $A=(-2,2)$. Giao với $B=\{-2,-1,0,1\}$ chỉ lấy các số nguyên nằm trong $(-2,2)$: $A\cap B=\{-1,0,1\}$. Hai đầu mút $-2$ và $2$ không thuộc $A$ nên không lấy.


\end{dapan}
\vspace{0.3cm}
\begin{dapan}[3 — Đáp án C]

\textbf{Giải.} Chia trục số theo hai mốc $-1$ và $2$. Với $x<-1$: $|x-2|+|x+1|=(2-x)+(-1-x)=1-2x<5 \iff x>-2$. Với $-1\le x\le2$: tổng luôn bằng $3<5$, nhận cả đoạn. Với $x>2$: $(x-2)+(x+1)=2x-1<5 \iff x<3$. Hợp lại được $(-2,3)$.


\end{dapan}
\vspace{0.3cm}
\begin{dapan}[4 — Đáp án A]

\textbf{Giải.} $x^2-(a+1)x+a=(x-1)(x-a)\le0$ nên $A$ là đoạn giữa $1$ và $a$ (hoặc $A=\{1\}$ khi $a=1$). Còn $B=(1,3)$, là khoảng mở nên không chứa $1$. Mọi phương án của $A$ đều chứa $1$, do đó $A\subseteq B$ không thể xảy ra: không có $a$ nào thỏa mãn.


\end{dapan}
\vspace{0.3cm}
\begin{dapan}[5 — Đáp án A]

\textbf{Giải.} Cauchy cho hai số dương: $x\cdot 2y\le\left(\dfrac{x+2y}{2}\right)^2=4$, suy ra $xy\le2$. Dấu bằng khi $x=2y=2$, tức $x=2$, $y=1$.

\section{PHẦN II — HÀM SỐ}


\end{dapan}
\vspace{0.3cm}
\begin{dapan}[6 — Đáp án B]

\textbf{Giải.} Điều kiện $1-x^2\ge0$, tức $x\in[-1,1]$. Phương trình $x=\sqrt{1-x^2}$ bình phương hai vế cho $x^2=1-x^2 \iff x^2=\dfrac{1}{2} \iff x=\pm\dfrac{\sqrt{2}}{2}$. Nghiệm $x=-\dfrac{\sqrt{2}}{2}$ bị loại vì vế trái âm còn vế phải không âm. Vậy có đúng một nghiệm.


\end{dapan}
\vspace{0.3cm}
\begin{dapan}[7 — Đáp án C]

\textbf{Giải.} Mẫu khác $0$: $x-1\ne0 \iff x\ne1$. Tập xác định là $\{x\mid x\ne1\}$.


\end{dapan}
\vspace{0.3cm}
\begin{dapan}[8 — Đáp án A]

\textbf{Giải.} Điều kiện $x^2>0 \iff x\ne0$, tập xác định đối xứng qua gốc. Và $f(-x)=\ln((-x)^2)=\ln(x^2)=f(x)$. Hàm chẵn.


\end{dapan}
\vspace{0.3cm}
\begin{dapan}[9 — Đáp án C]

\textbf{Giải.} Mệnh đề C sai vì cơ số quyết định chiều biến thiên: $y=a^x$ đồng biến khi $a>1$ và nghịch biến khi $0<a<1$. Ba mệnh đề còn lại đúng với mọi cơ số hợp lệ.


\end{dapan}
\vspace{0.3cm}
\begin{dapan}[10 — Đáp án C]

\textbf{Giải.} Từ $y=\dfrac{1}{2}x+3$ suy ra $x=2(y-3)=2y-6$. Đổi tên biến: $y=2x-6$ với $x\in\mathbb{R}$.


\end{dapan}
\vspace{0.3cm}
\begin{dapan}[11 — Đáp án D]

\textbf{Giải.} $\cos(-\alpha)=\cos\alpha$ đúng vì cosin là hàm chẵn. Ba hệ thức còn lại sai dấu: $\sin(\pi+\alpha)=-\sin\alpha$, $\cos(\pi+\alpha)=-\cos\alpha$, $\sin(-\alpha)=-\sin\alpha$.

\section{PHẦN III — LƯỢNG GIÁC}


\end{dapan}
\vspace{0.3cm}
\begin{dapan}[12 — Đáp án D]

\textbf{Giải.} $\dfrac{5\pi}{6}$ thuộc góc phần tư thứ hai nên cosin âm. $\cos\dfrac{5\pi}{6}=-\cos\left(\pi-\dfrac{5\pi}{6}\right)=-\cos\dfrac{\pi}{6}=-\dfrac{\sqrt{3}}{2}$.


\end{dapan}
\vspace{0.3cm}
\begin{dapan}[13 — Đáp án B]

\textbf{Giải.} Chu kỳ của $y=A\cos(\omega x+\varphi)$ là $T=\dfrac{2\pi}{\lvert\omega\rvert}$. Với $\omega=4$: $T=\dfrac{2\pi}{4}=\dfrac{\pi}{2}$. Hệ số $3$ và pha $-\dfrac{\pi}{3}$ không ảnh hưởng chu kỳ.


\end{dapan}
\vspace{0.3cm}
\begin{dapan}[14 — Đáp án B]

\textbf{Giải.} $\alpha\in\left(\dfrac{\pi}{2},\pi\right)$ nên $\dfrac{\alpha}{2}\in\left(\dfrac{\pi}{4},\dfrac{\pi}{2}\right)$, góc phần tư thứ nhất, sin dương. Công thức hạ bậc: $\sin\dfrac{\alpha}{2}=\sqrt{\dfrac{1-\cos\alpha}{2}}=\sqrt{\dfrac{1+\frac{1}{2}}{2}}=\sqrt{\dfrac{3}{4}}=\dfrac{\sqrt{3}}{2}$.


\end{dapan}
\vspace{0.3cm}
\begin{dapan}[15 — Đáp án C]

\textbf{Giải.} $\alpha\in\left(-\dfrac{\pi}{2},0\right)$ nên $\dfrac{\alpha}{2}\in\left(-\dfrac{\pi}{4},0\right)$, góc phần tư thứ tư, tang âm. Công thức: $\tan\dfrac{\alpha}{2}=-\sqrt{\dfrac{1-\cos\alpha}{1+\cos\alpha}}=-\sqrt{\dfrac{1-\frac{1}{3}}{1+\frac{1}{3}}}=-\sqrt{\dfrac{1}{2}}=-\dfrac{\sqrt{2}}{2}$.


\end{dapan}
\vspace{0.3cm}
\begin{dapan}[16 — Đáp án D]

\textbf{Giải.} Công thức góc phụ: $\cos\left(\dfrac{\pi}{2}+\alpha\right)=-\sin\alpha=-\dfrac{1}{3}$.


\end{dapan}
\vspace{0.3cm}
\begin{dapan}[17 — Đáp án B]

\textbf{Giải.} $\dfrac{\pi}{12}=\dfrac{\pi}{4}-\dfrac{\pi}{6}$, vậy $\cos\dfrac{\pi}{12}=\cos\dfrac{\pi}{4}\cos\dfrac{\pi}{6}+\sin\dfrac{\pi}{4}\sin\dfrac{\pi}{6}=\dfrac{\sqrt{2}}{2}\cdot\dfrac{\sqrt{3}}{2}+\dfrac{\sqrt{2}}{2}\cdot\dfrac{1}{2}=\dfrac{\sqrt{6}+\sqrt{2}}{4}$.


\end{dapan}
\vspace{0.3cm}
\begin{dapan}[18 — Đáp án B]

\textbf{Giải.} $\alpha$ góc phần tư thứ nhất nên $\cos\alpha>0$, suy ra $\cos\alpha=\sqrt{1-\dfrac{9}{25}}=\dfrac{4}{5}$. Công thức nhân đôi: $\sin 2\alpha=2\sin\alpha\cos\alpha=2\cdot\dfrac{3}{5}\cdot\dfrac{4}{5}=\dfrac{24}{25}$.


\end{dapan}
\vspace{0.3cm}
\begin{dapan}[19 — Đáp án B]

\textbf{Giải.} $\alpha\in\left(\pi,\dfrac{3\pi}{2}\right)$ nên $\dfrac{\alpha}{2}\in\left(\dfrac{\pi}{2},\dfrac{3\pi}{4}\right)$, cosin âm. Từ $\sin\alpha=-\dfrac{12}{13}$ suy ra $\cos\alpha=-\dfrac{5}{13}$. Vậy $\cos\dfrac{\alpha}{2}=-\sqrt{\dfrac{1-\frac{5}{13}}{2}}=-\sqrt{\dfrac{4}{13}}=-\dfrac{2\sqrt{13}}{13}$.


\end{dapan}
\vspace{0.3cm}
\begin{dapan}[20 — Đáp án A]

\textbf{Giải.} Theo định nghĩa, $\tan x=\dfrac{\sin x}{\cos x}$ với điều kiện $\cos x\ne0$. Các phương án còn lại nhầm với $\sec x=\dfrac{1}{\cos x}$ và $\cot x=\dfrac{\cos x}{\sin x}$.


\end{dapan}
\vspace{0.3cm}
\begin{dapan}[21 — Đáp án D]

\textbf{Giải.} Hàm $y=\sin x$ đồng biến trên $\left[-\dfrac{\pi}{2}+2k\pi,\dfrac{\pi}{2}+2k\pi\right]$. Với $k=0$ được $\left[-\dfrac{\pi}{2},\dfrac{\pi}{2}\right]$. Trên $[0,\pi]$ hàm tăng rồi giảm, còn trên toàn trục số hàm không đơn điệu.

\section{PHẦN IV — DÃY SỐ}


\end{dapan}
\vspace{0.3cm}
\begin{dapan}[22 — Đáp án C]

\textbf{Giải.} Gọi $G$ là trung bình nhân: $G^2=3\cdot12=36$, suy ra $G=\pm6$.


\end{dapan}
\vspace{0.3cm}
\begin{dapan}[23 — Đáp án D]

\textbf{Giải.} $a_n=a_1+(n-1)d$, vậy $a_{100}=2+99\cdot3=2+297=299$.


\end{dapan}
\vspace{0.3cm}
\begin{dapan}[24 — Đáp án C]

\textbf{Giải.} $a_4=a_1q^3 \Rightarrow 32=4q^3 \Rightarrow q^3=8 \Rightarrow q=2$. Do đó $a_3=a_1q^2=4\cdot4=16$.


\end{dapan}
\vspace{0.3cm}
\begin{dapan}[25 — Đáp án A]

\textbf{Giải.} Từ $\dfrac{1}{a_{n+1}}-\dfrac{1}{a_n}=1$, dãy $\left\{\dfrac{1}{a_n}\right\}$ là cấp số cộng với số hạng đầu $\dfrac{1}{a_1}=2$ và công sai $1$. Vậy $\dfrac{1}{a_n}=2+(n-1)=n+1$, suy ra $a_n=\dfrac{1}{n+1}$ và $a_6=\dfrac{1}{7}$.


\end{dapan}
\vspace{0.3cm}
\begin{dapan}[26 — Đáp án D]

\textbf{Giải.} Dùng tính chất ghép cặp của cấp số cộng. Ta có $a_2+a_6=2a_4$ và $a_6+a_{10}+a_{14}=3a_{10}$. Điều kiện đã cho thành $3\cdot2a_4+2\cdot3a_{10}=24$, tức $6(a_4+a_{10})=24$, suy ra $a_4+a_{10}=4$. Mặt khác $a_4+a_{10}=2a_7$ nên $a_7=2$. Cuối cùng $S_{13}=\dfrac{13(a_1+a_{13})}{2}=\dfrac{13\cdot2a_7}{2}=13a_7=26$.


\end{dapan}
\vspace{0.3cm}
\begin{dapan}[27 — Đáp án B]

\textbf{Giải.} Dấu đan xen bắt đầu bằng âm nên tử số chứa $(-1)^n$. Mẫu số là $5,7,9,11,\ldots$, tức $2n+3$. Vậy $a_n=\dfrac{(-1)^n}{2n+3}$. Thử $n=1$: $-\dfrac{1}{5}$, đúng.

\section{PHẦN V — ÔN TẬP TÍCH LŨY}


\end{dapan}
\vspace{0.3cm}
\begin{dapan}[28 — Đáp án D]

\textbf{Giải.} $\displaystyle\int_0^1 (3x^2+1)\,dx=\left[x^3+x\right]_0^1=1+1=2$.


\end{dapan}
\vspace{0.3cm}
\begin{dapan}[29 — Đáp án A]

\textbf{Giải.} Nguyên hàm $f(x)=x^3-x+C$. Điều kiện $f(1)=1-1+C=0$ cho $C=0$. Vậy $f(x)=x^3-x$.


\end{dapan}
\vspace{0.3cm}
\begin{dapan}[30 — Đáp án A]

\textbf{Giải.} Đưa hằng số ra ngoài: $\displaystyle\int\dfrac{1}{3x}\,dx=\dfrac{1}{3}\int\dfrac{1}{x}\,dx=\dfrac{1}{3}\ln\lvert x\rvert+C$.


\end{dapan}
\vspace{0.3cm}
\begin{dapan}[31 — Đáp án A]

\textbf{Giải.} Hàm hợp với $u=2x+3$: $\left[\sin(2x+3)\right]'=\cos(2x+3)\cdot2=2\cos(2x+3)$.


\end{dapan}
\vspace{0.3cm}
\begin{dapan}[32 — Đáp án A]

\textbf{Giải.} $y'=3x^2+3$, tại $x=0$ có hệ số góc $y'(0)=3$. Tiếp tuyến qua $P(0,1)$: $y=3x+1$, viết lại thành $3x-y+1=0$.


\end{dapan}
\vspace{0.3cm}
\begin{dapan}[33 — Đáp án C]

\textbf{Giải.} Viết $y=x^{1/2}$, đạo hàm $y'=\dfrac{1}{2}x^{-1/2}=\dfrac{1}{2\sqrt{x}}$.


\end{dapan}
\vspace{0.3cm}
\begin{dapan}[34 — Đáp án B]

\textbf{Giải.} Với $0<a<1$, hàm $y=a^x$ nghịch biến trên $\mathbb{R}$, tức là hàm giảm. Hàm mũ không có tính chẵn lẻ.


\end{dapan}
\vspace{0.3cm}
\begin{dapan}[35 — Đáp án D]

\textbf{Giải.} Với $x<0$: $x+2<4 \iff x<2$, nhận cả $x<0$. Với $x\ge0$: $x^2<4 \iff 0\le x<2$. Hợp lại được $x<2$, tức $(-\infty,2)$.


\end{dapan}
\vspace{0.3cm}
\begin{dapan}[36 — Đáp án C]

\textbf{Giải.} Đạo hàm hàm hợp: $y'=4\cdot2(2-x+3x^2)\cdot(-1+6x)=8(2-x+3x^2)(6x-1)$.


\end{dapan}
\vspace{0.3cm}
\begin{dapan}[37 — Đáp án A]

\textbf{Giải.} $f'(x)=\cos x$, nên $f'(0)=\cos0=1$.


\end{dapan}
\vspace{0.3cm}
\begin{dapan}[38 — Đáp án C]

\textbf{Giải.} $f(-x)=-3x+1$, so sánh: $f(-x)\ne f(x)$ (vì $1\ne-1$) và $f(-x)\ne-f(x)$ (vì $-3x+1\ne-3x-1$). Hàm không chẵn không lẻ.


\end{dapan}
\vspace{0.3cm}
\begin{dapan}[39 — Đáp án B]

\textbf{Giải.} Hàm lũy thừa có dạng $y=x^\alpha$. Chỉ $y=x^{-1}$ đúng dạng; $y=2x^2$ và $y=3x$ có thêm hệ số, còn $y=x^3-1$ có thêm hằng số.


\end{dapan}
\vspace{0.3cm}
\begin{dapan}[40 — Đáp án D]

\textbf{Giải.} Cơ số $5>1$ nên bất phương trình tương đương $2x+1\ge1 \iff x\ge0$. Kết hợp điều kiện xác định $2x+1>0 \iff x>-\dfrac{1}{2}$, nghiệm là $x\ge0$.


\end{dapan}
\vspace{0.3cm}
\begin{dapan}[41 — Đáp án C]

\textbf{Giải.} $a+2b=\lg4+2\lg5=\lg4+\lg25=\lg100=2$.


\end{dapan}
\vspace{0.3cm}
\begin{dapan}[42 — Đáp án D]

\textbf{Giải.} Hàm $\log_a u$ bằng $0$ khi $u=1$, bất kể cơ số. Vậy $x-3=1 \iff x=4$, và $f(4)=2+0=2$. Điểm cố định là $(4,2)$.


\end{dapan}
\vspace{0.3cm}
\begin{dapan}[43 — Đáp án D]

\textbf{Giải.} Cơ số $\dfrac{1}{2}<1$ nên hàm nghịch biến trên $\mathbb{R}$, tức là hàm giảm. Đồ thị nằm hoàn toàn phía trên trục hoành nên không đối xứng qua trục $x$ cũng không qua gốc.


\end{dapan}
\vspace{0.3cm}
\begin{dapan}[44 — Đáp án A]

\textbf{Giải.} Số sách ít nhất phải là $70+50-30=90$ cuốn, nhỏ hơn $100$ nên dữ liệu không mâu thuẫn. Hai nhóm không loại trừ nhau, vì $30$ cuốn vừa là tiểu thuyết vừa là khoa học.


\end{dapan}
\vspace{0.3cm}
\begin{dapan}[45 — Đáp án B]

\textbf{Giải.} Số người biết ít nhất một môn: $18+15-5=28$. Nhóm $30$ người nên số người không biết môn nào là $30-28=2$.


\end{dapan}
\vspace{0.3cm}
\begin{dapan}[46 — Đáp án B]

\textbf{Giải.} Với $x<0$, đặt $t=-x>0$ thì $f(x)=-\left(t+\dfrac{1}{t}\right)\le-2$ theo Cauchy. Dấu bằng khi $t=1$, tức $x=-1$. Giá trị lớn nhất là $-2$.


\end{dapan}
\vspace{0.3cm}
\begin{dapan}[47 — Đáp án D]

\textbf{Giải.} Mẫu khác $0$: $x\ne\dfrac{1}{2}$. Bất phương trình $\lvert 2x-1\rvert\le1 \iff -1\le2x-1\le1 \iff 0\le x\le1$. Loại điểm $x=\dfrac{1}{2}$ được $\left[0,\dfrac{1}{2}\right)\cup\left(\dfrac{1}{2},1\right]$.


\end{dapan}
\vspace{0.3cm}
\begin{dapan}[48 — Đáp án A]

\textbf{Giải.} Tam thức $(x-a)(x-b)$ âm trong khoảng giữa hai nghiệm, nên $\{a,b\}=\{5,9\}$. Tổng $a+b=14$.
\end{dapan}
\vspace{0.3cm}

\vspace{0.8cm}
\begin{center}
{\footnotesize\color{cae_blue} Tài liệu nội bộ, CAE SHANGHAI}
\end{center}

\end{document}
